% ===== DEVOPS =====
\section{DevOps: Infraestrutura, Automação e Entrega Contínua}

DevOps \cite{wang2023} integra desenvolvimento, operação e observabilidade para que a plataforma evolua com segurança. No SmartHive, a estratégia adotada prioriza a reprodutibilidade do ambiente por meio de containers Docker, permitindo executar frontend e backend de forma padronizada em ambiente local, acadêmico e de homologação. A aplicação web foi conteinerizada em dois serviços principais: a API FastAPI em Python e o frontend React/Vite servido por Nginx.

\subsection{Pipeline CI/CD}

O pipeline de integração e entrega contínua apoia validação automática, construção das imagens Docker e preparação dos ambientes:

\begin{figure}[H]
\centering
\begin{tikzpicture}[
    node distance=1.5cm,
    every node/.style={rectangle, draw, rounded corners,
                       fill=secondarycolor!20,
                       text width=2cm,
                       align=center,
                       font=\small,
                       minimum height=1cm}
]
\node (code) {Código};
\node (build) [right of=code, xshift=1cm] {Build Docker};
\node (test) [right of=build, xshift=1cm] {Validação};
\node (deploy) [right of=test, xshift=1cm] {Deploy};
\node (monitor) [right of=deploy, xshift=1cm] {Monitorar};

\draw[->, thick] (code) -- (build);
\draw[->, thick] (build) -- (test);
\draw[->, thick] (test) -- (deploy);
\draw[->, thick] (deploy) -- (monitor);
\draw[->, thick, dashed] (monitor) to[bend right=45] (code);
\end{tikzpicture}
\caption{Pipeline CI/CD ajustado para a conteinerização Docker}
\label{fig:cicd}
\end{figure}

\subsection{Exemplo de Configuração CI/CD}

Exemplo de workflow mínimo em GitHub Actions para validar a aplicação conteinerizada:

\begin{verbatim}
name: ci

on:
  push:
    branches: [ main ]
  pull_request:

jobs:
  validate:
    runs-on: ubuntu-latest
    defaults:
      run:
        working-directory: web-smarthive
    steps:
      - uses: actions/checkout@v4
      - uses: actions/setup-node@v4
        with:
          node-version: 20
      - uses: actions/setup-python@v5
        with:
          python-version: "3.12"
      - run: npm ci
        working-directory: web-smarthive/frontend
      - run: npm run build
        working-directory: web-smarthive/frontend
      - run: pip install -r requirements.txt
        working-directory: web-smarthive/backend
      - run: docker compose build
      - run: docker compose config
\end{verbatim}

\subsection{Infraestrutura como Código}

\begin{table}[H]
\centering
\caption{Ferramentas DevOps Utilizadas}
\label{tab:devops-tools}
\begin{tabular}{|l|l|p{6cm}|}
\hline
\textbf{Categoria} & \textbf{Ferramenta} & \textbf{Finalidade} \\
\hline
Controle de Versão & Git + GitHub & Versionamento centralizado na branch \texttt{main} \\
CI/CD & GitHub Actions & Validação de build, dependências e imagens Docker \\
Containerização & Docker & Empacotamento do backend FastAPI e frontend React/Vite \\
Orquestração local & Docker Compose & Execução integrada dos serviços \texttt{backend} e \texttt{frontend} \\
Servidor web & Nginx & Entrega dos arquivos estáticos do frontend e proxy para \texttt{/api} \\
Banco e autenticação & Supabase & Persistência PostgreSQL, autenticação e storage compartilhado \\
Configuração & Variáveis de ambiente & Parametrização de portas, URLs, Supabase e modo de execução \\
\hline
\end{tabular}
\end{table}

\subsection{Infraestrutura do Sistema}

Esta seção documenta como o sistema funciona em diferentes ambientes e plataformas.

\subsubsection{Componentes principais}

\begin{table}[H]
\centering
\caption{Componentes de Infraestrutura}
\label{tab:infrastructure}
\begin{tabular}{|l|p{10cm}|}
\hline
\textbf{Componente} & \textbf{Descrição} \\
\hline
\textbf{Camada de Captura} &
\begin{itemize}[leftmargin=*, nosep]
    \item Aplicativo mobile e sensores como fontes de coleta em campo
    \item Envio de dados de apiários, colmeias, capturas e monitoramentos
    \item Sincronização com o mesmo banco Supabase utilizado pelo web
\end{itemize} \\
\hline
\textbf{Backend e API} &
\begin{itemize}[leftmargin=*, nosep]
    \item Container \texttt{smarthive-backend:local} baseado em \texttt{python:3.12-slim}
    \item API FastAPI executada com \texttt{uvicorn} na porta \texttt{8000}
    \item \texttt{HEALTHCHECK} em \texttt{/api/health} para validar disponibilidade do serviço
\end{itemize} \\
\hline
\textbf{Persistência e Analytics} &
\begin{itemize}[leftmargin=*, nosep]
    \item Supabase/PostgreSQL para autenticação, perfis, apiários, colmeias e eventos
    \item Volumes Docker para dados locais e uploads quando o modo Supabase não estiver ativo
    \item Variáveis \texttt{DATABASE\_URL}, \texttt{SUPABASE\_URL} e chaves do Supabase para integração
\end{itemize} \\
\hline
\textbf{Aplicações Cliente} &
\begin{itemize}[leftmargin=*, nosep]
    \item Container \texttt{smarthive-frontend:local} com build Node 20 e entrega por Nginx
    \item Frontend React/Vite servido na porta configurável \texttt{FRONTEND\_PORT}
    \item Nginx encaminhando \texttt{/api} e \texttt{/uploads} para o backend no Docker Compose
\end{itemize} \\
\hline
\textbf{Segurança e Observabilidade} &
\begin{itemize}[leftmargin=*, nosep]
    \item Autenticação centralizada no Supabase Auth para web e mobile
    \item CORS parametrizado por variáveis de ambiente
    \item Healthchecks nos containers para facilitar reinício e diagnóstico operacional
\end{itemize} \\
\hline
\end{tabular}
\end{table}

\subsubsection{Diagrama conceitual de dados}

\begin{figure}[H]
\centering
\includegraphics[width=\textwidth,height=0.72\textheight,keepaspectratio]{\ArchitectureDiagramPath}
\caption{Fluxo conceitual de captura, processamento e consumo dos dados}
\label{fig:infrastructure-diagram}
\end{figure}

\textbf{Leitura do fluxo ilustrado:}

\begin{itemize}
    \item \textbf{Captura:} reúne dados do aplicativo mobile, sensores e entradas manuais
    \item \textbf{Processamento:} organiza requisições na API FastAPI executada em container Docker
    \item \textbf{Persistência:} centraliza autenticação e dados no Supabase/PostgreSQL
    \item \textbf{Consumo:} entrega dashboard web via Nginx e aplicativo mobile usando a mesma base
\end{itemize}

\subsection{Arquitetura de Deploy}

Uma arquitetura de deploy enxuta para o projeto pode ser organizada da seguinte forma:

\begin{figure}[H]
\centering
\begin{tikzpicture}[
    node distance=2cm,
    component/.style={rectangle, draw, rounded corners,
                      fill=secondarycolor!20,
                      text width=2.5cm,
                      align=center,
                      minimum height=1.2cm}
]
\node[component] (nginx) {Frontend\\Nginx};
\node[component, below left of=nginx] (api) {Backend\\FastAPI};
\node[component, below right of=nginx] (mobile) {Aplicativo\\Mobile};
\node[component, below of=api, xshift=1.5cm] (db) {Supabase\\PostgreSQL/Auth};
\node[component, below of=db] (volumes) {Volumes\\Uploads locais};

\draw[<->, thick] (nginx) -- (api);
\draw[<->, thick] (mobile) -- (api);
\draw[<->, thick] (api) -- (db);
\draw[<->, thick] (mobile) -- (db);
\draw[<->, thick] (api) -- (volumes);
\end{tikzpicture}
\caption{Arquitetura de deploy conteinerizada com Docker Compose}
\label{fig:deployment}
\end{figure}

\subsection{Métricas e SLA}

Para acompanhamento operacional, algumas métricas relevantes são:

\begin{itemize}
    \item \textbf{Disponibilidade do painel:} 99.5\% de uptime mensal
    \item \textbf{Latência média da API:} abaixo de 500ms para consultas principais
    \item \textbf{Tempo de subida local:} containers saudáveis após \texttt{docker compose up --build}
    \item \textbf{Saúde dos serviços:} healthchecks ativos para frontend e backend
    \item \textbf{Integridade dos dados:} persistência centralizada no Supabase e volumes para arquivos locais
    \item \textbf{Lead time de mudança:} até 1 dia para ajustes de baixa complexidade
    \item \textbf{MTTR:} até 2 horas para restauração de falhas operacionais
\end{itemize}
